\bookchapter{Bab 0: Hai, sebelum memulai}

Sebelum menjelajah secara resmi, mari kita bahas beberapa aturan dasar dulu, untuk memastikan Anda memahami semua yang dibahas nanti. Jangan terburu-buru, menyiapkan pisau dengan baik tidak akan menghambat pekerjaan menebang kayu.

\booksection{Apa yang bisa dibantu oleh buku ini? (dan yang tidak bisa)}

Buku ini bisa membantu Anda memahami secara sistematis cara menggunakan ponsel, memahami alasan di balik fitur-fiturnya, dan baik Anda pengguna Android maupun Apple, sebagian besar konsepnya berlaku umum. Buku ini akan mengajak Anda berlatih langsung, meningkatkan Anda dari “bisa menggunakan” menjadi “memahami penggunaan”.

Buku ini bukan manual untuk ponsel tertentu, juga bukan berita teknologi tentang produk terbaru, apalagi panduan untuk membuat Anda kecanduan ponsel. Sejujurnya, jika yang Anda harapkan adalah “cara menonton video di ponsel dengan lebih nikmat”, Anda mungkin akan kecewa—kami lebih peduli pada cara membuat ponsel benar-benar melayani Anda, bukan sebaliknya.

\booksection{Penjelasan “kode rahasia” kami}

Untuk menjelaskan dengan lebih jelas, buku ini menggunakan beberapa penanda kecil, penanda ini seperti kode rahasia di antara kita:

\begin{itemize}
  \item \textbf{[Nama Tombol]}: Mewakili tombol atau nama aplikasi spesifik di layar, misalnya [Pengaturan], [OK].
  \item \textbf{>}: Mewakili urutan operasi, misalnya "klik [Pengaturan] > [Suara]", berarti klik Pengaturan dulu, lalu klik Suara.
  \item \textbf{Tebal}: Sangat penting! Mungkin konsep kunci atau peringatan keamanan, harap perhatikan.
  \item \emph{Miring}: Beberapa pemikiran atau perumpamaan tambahan, untuk membantu Anda memahami dengan lebih baik.
\end{itemize}

\booksection{Android dan iOS, kami bahas keduanya}

Tebak apa, beberapa hal namanya berbeda, tapi fungsinya sama. Misalnya: aplikasi (Android) dan App (iOS), pengaturan (keduanya disebut pengaturan), manajer file (Android) dan File (iOS). Yang penting Anda tahu itu hal yang sama, jangan bingung dengan istilahnya.

Selain itu, di kubu Android ada juga sistem HarmonyOS milik Huawei, yang memiliki karakteristiknya sendiri, akan kita bahas khusus nanti. Singkatnya, apa pun ponsel yang Anda gunakan, buku ini bisa membantu Anda.

\booksection{Kapasitas ponsel: Mengapa membeli 128G, tapi yang tersedia kurang dari 110G?}

Unit kapasitas yang diketahui semua: 1TB = 1024GB, 1GB = 1024MB, 1MB = 1024KB... tapi ada jebakan di sini: produsen hard disk biasanya menghitung dengan basis 1000, sedangkan sistem menghitung dengan basis 1024, jadi hard disk “128G” yang Anda beli, kapasitas aktualnya akan “menyusut”.

Tapi alasan utama kapasitas ponsel Anda menyusut ada tiga:

\begin{enumerate}
  \item \textbf{Sistem itu sendiri menempati ruang}: Sistem operasi seperti dinding penyangga rumah, harus ada. Sistem sekarang sering menempati 10-20GB, ini sangat normal.
  \item \textbf{Aplikasi pre-installed dari produsen}: Aplikasi yang sudah ada di ponsel baru juga memakan ruang, beberapa bisa di-uninstal, beberapa tidak.
  \item \textbf{Kehilangan partisi dan format}: Alasan teknis, seperti lapisan kain tas juga memakan volume.
\end{enumerate}

\begin{tcbinfo}{Saran kecil saat membeli ponsel}
Saat ini ponsel 128GB pada dasarnya tidak cukup, mulai dari 256GB lebih aman. Jika Anda suka fotografi atau bermain game, disarankan langsung ke 512GB atau lebih besar. Ingat, untuk ruang penyimpanan, lebih baik berlebih daripada kurang.

\end{tcbinfo}

\begin{tcbpractice}{Mari praktik: Lihat "KTP" ponsel Anda}
Buka [Pengaturan], temukan tentang ponsel (atau opsi serupa), catat: model ponsel, versi sistem, total ruang penyimpanan dan ruang tersisa. Ini adalah langkah pertama untuk mengenalnya. Sekalian cek kesehatan baterai, supaya ada gambaran.

\end{tcbpractice}

\booksection{Ponsel Anda adalah “gudang aset digital”}

Sebelum memulai pengetahuan spesifik, mari kita bangun pemahaman inti:\textbf{ponsel bukan lagi sekadar alat komunikasi, ia adalah gudang aset digital terpenting kita}. Kontak, riwayat chat, foto keluarga, dokumen kerja, aplikasi bank, kode pembayaran, data kesehatan, dokumen elektronik... data-data ini digabungkan, nilainya sering melebihi ponsel itu sendiri.

\begin{tcbinfo}{Mengapa ini penting?}
Jika ponsel hilang atau rusak, yang Anda hilang bukan hanya perangkat keras senilai beberapa ribu yuan,\textbf{tapi lebih mungkin data yang tidak bisa diperbarui yang terkumpul bertahun-tahun}—foto tumbuh kembang anak, suara kerabat yang telah meninggal, kontrak penting yang tidak dicadangkan. Kebiasaan mengatur dan melindungi secara ilmiah bisa melindungi Anda dari tiga dimensi:

\begin{itemize}
  \item \textbf{Meningkatkan efisiensi operasi}: Klasifikasi file yang teratur membuat pencarian sistem lebih cepat, ponsel tidak lag dalam waktu lama.
  \item \textbf{Memperpanjang usia pakai}: Manajemen penyimpanan yang wajar mengurangi beban baca-tulis chip, baterai dan penyimpanan lebih tahan lama.
  \item \textbf{Menjamin keamanan data}: Manajemen berlapis + pencadangan berkala + pembatasan izin, membuat "bencana data" jauh dari Anda.
\end{itemize}

\end{tcbinfo}

\booksubsection{Filosofi inti kami: “Pencegahan utama, kombinasi pencegahan dan penanganan”}

Sama seperti menjaga kesehatan tidak bisa menunggu sakit baru ke dokter,\textbf{manajemen ponsel juga harus dimulai dari “kebiasaan harian”}. Yang ditekankan berulang-ulang dalam buku ini bukan teknik tingkat lanjut tertentu, melainkan prinsip-prinsip sederhana ini:

\begin{itemize}
  \item \textbf{Pencegahan utama}: Instal aplikasi hanya dari toko resmi, hilangkan informasi sensitif foto sebelum dikirim, jangan gunakan kata sandi yang sama berulang—"tindakan kecil" ini biayanya sangat rendah, tapi bisa menangkal sebagian besar risiko.
  \item \textbf{Pemeriksaan berkala}: Atur folder unduhan sebulan sekali, audit izin sekali per kuartal, cadangan penuh sekali per enam bulan—ritme tetap jauh lebih efektif daripada pemadam kebakaran sesaat.
  \item \textbf{Sisakan jalan mundur}: Aktifkan sinkronisasi cloud, catat akun dan kata sandi dengan baik, tulis sebuah "daftar warisan digital"—pastikan ketika kejadian tak terduga terjadi, Anda dan keluarga bisa menghadapinya dengan tenang.
\end{itemize}
